\chapter{CONCLUSION}
\label{chap:conclusions}


After an extensive set of tests, we can finally draw some general and specific conclusions. First, spectral bias is undeniably present in all architectures, as a general behavior of training dynamics. The data supervision helps approximate the low and mid spectral bands better than what physics-informed showed. On the other hand, physics-informed brought more stability on the time axis since it helps to mitigate the Gibbs artifact coming from Fourier transforms on non-periodic time. The sections below read that result family by family, and then set out the boundaries of the study and the work that should follow it.

% ===============================================================
\section{Assessment of Each Method Family}
\label{sec:c_verdict}

The four architectures split into two families by how one trained network relates to the parameter axis, and reading them family by family also answers the questions posed in Chapter~\ref{chap:intro}.

\paragraph{Neural Networks: MLP and PINN} Recall that this family of methods aims to learn the solution for only one parameter. The MLP (a ``pure data PINN'') shows spectral bias in its plain form, fitting spectra as training goes. The pure physics PINN contrasts with it in the sense that it does not fit well mid and high bands of the spectrum: PINNs do not seem to approximate wave profile tails, meaning the high frequency content. Thus, the combination data + physics seems to be the best, gathering both strengths, the physics-informed part of the loss working as a regularization term a priori of the solution.

\paragraph{Neural Operators: FNO and PINO}

In this family, we have one architecture capable of showing how the solution looks at any model parameter close to what it saw in the dataset. Here we have similar results. The FNO, the pure data supervised architecture, simply learns to fit and, while still showing spectral bias, has its main weakness in the Gibbs phenomenon anomaly from taking Fourier transforms in the time axis. This introduces error shootouts at the temporal boundaries. Meanwhile, the pure physics PINO presents the same multiparameter solution, but shows the same behavior as the pure physics PINN, meaning they cannot approximate high frequency details of the solution or, in other words, their ripples. The best method is, again, the mix of data + physics, and here we can observe that the Gibbs anomaly in the FNO is suppressed in this supervised PINO, ultimately presenting this work's best architecture, which fits the data and respects the model's physics. Also, discretization invariance holds for the two architectures, with only one caveat: increasing resolution yields slightly lower quality results, but the architectures remain advantageous as surrogates.


\paragraph{Spectral Bias Study}

After we gathered, to the best of our knowledge, theoretical findings that explain or expose spectral bias, we condensed them in this work and designed the experiments to verify its plausibility. For all we can say, while the simple network struggled to learn the target, the kernel stayed close to its initial setting and thus our theory predicts the loss behavior well on each NTK kernel eigenvector direction (or mode, for our problem). The network size also shows us that, really, increasing the size of the MLP indeed helps the kernel to change less, allowing for the good estimates presented, such as loss evolution by mode and iteration a priori estimate. As long as the kernel does not drift (drifting is a good thing, it means the network is learning), our predictions are good.

The cost accounting of Table~\ref{tab:training} completes the comparison. A pointwise network is the cheaper object to train and the more expensive one to use, since a new parameter means a new training run, while an operator pays its training cost once and answers the rest of the sweep in forward passes. Both physics-informed variants cost more to train than their data-driven counterparts, since every iteration differentiates the network to form the residual.

% ===============================================================
\section{What This Study Cannot Claim}
\label{sec:c_limits}

Several boundaries fence in the results above, and they should be weighed before the conclusions are carried anywhere else.

Every experiment used one initial condition, the unit solitary profile of Table~\ref{tab:m_fixed}, and held the two coefficients equal. The operators therefore learned a one-parameter family and not the four-parameter Boussinesq family classified by \textcite{bona2002} and \textcite{martins2023}, so the generalization reported here is generalization along a single coordinate. Decoupling $\alpha$ from $\beta$, or varying the initial data, is a different potential challenge to be explored.

The temporal advantage measured for the PINO is entangled with an implementation choice. Our operators transform along the time axis as well as the spatial one, and the terminal-time artifact of Section~\ref{sec:r_gibbs} follows from that assumption of temporal periodicity. An operator that marched in time, as \textcite{li2021} suggest, would not produce the artifact at all, so part of what the residual is credited with here is the repair of a defect the architecture introduced.

Capacity and budget were fixed rather than tuned. Both operators keep $16$ Fourier modes per axis, every trained method ran modest $15{,}000$ Adam iterations at a single learning rate, and the whole pipeline ran from one seed on one GPU.

The reference solution is itself a numerical solution, accurate to the order of the pseudospectral scheme on the grid it was solved on and not exact, which matters at the high wavenumbers, where it carries the largest share of its own error and where the trained methods are weakest.

The minimal experiment is a model of the mechanism and not a model of the trained methods. It fixes time, fits a single profile by plain gradient descent, and works with a kernel small enough to be eigendecomposed at every checkpoint. Not much of this analysis transfers easily to PINNs and NOs.

% ===============================================================
\section{Considerations and Directions Forward}
\label{sec:c_future}

For directions: adaptive weighting of the loss terms \parencite{wang2021understanding}, the Neural Tangent Kernel weighting derived for PINNs by \textcite{wang2022ntk}, and the second-order and spectrum-aware optimizers of \textcite{khodakarami2026} are possible near-future steps to consider, evaluating possible mitigations for spectral bias.

An operator that marches in time removes the temporal periodicity assumption by construction, which would isolate how much of the late-time steadiness of the PINO comes from the physics term itself. A related question is whether the residual keeps helping once the training data are made scarce, the regime the physics-informed formulation was proposed for and the one this study, with a reference solution available at every point of the grid, did not test. Maybe we could try some real problems with observations that include errors and whatnot.

The decoupling of the nonlinearity from the dispersion, varying the initial data, and extending to the two-dimensional system would show whether the ordering found here survives a solution operator of higher dimension, and the analysis of \textcite{krishnapriyan2021characterizing} gives reason to expect the physics-informed losses to grow harder to optimize as that happens.

We can also explore other architectures like DeepONet \parencite{lu2021deeponet} and hybrid schemes that hand the high-frequency residual back to a classical corrector, which would test whether the complementary role identified here can be built into a single solver.

% ===============================================================

At this time, we can accept spectral bias as a behavior we can choose to mitigate, and we recommend doing so, since the literature presents simple and effective ways. It should be taken into account and acknowledged, if the reader is aiming to implement these architectures. This also shows that ML methods do not ultimately come to take the place of numerical methods. Key Fluid Dynamics simulations in the industry will not stop using classical-based solvers for ML methods, as these should actually work together. In the end, a trained operator answers a parameter test sweep for parametric physical models in a single forward pass, for a fraction of what repeated solves would cost. Also, unlike classical methods, theoretical analyses for ML methods are much harder to come by and not yet fully consolidated.

All of this considered, we finish this work noting that this is a novel field and most certainly the new era of Fluid Dynamics where ML models promise fast evaluation and new predictions. Of course this is not tied to only neural networks and neural operators, given the vastness of the field.
